\begin{figure}[ht]
  \centering
  % NIENTE \resizebox: larghezze calibrate per impedire il downscaling ottico
  \begin{tikzpicture}[
      % Stili Gold Standard
      node distance=4ex and 2.5em,
      %
      cls/.style={draw, rounded corners=1pt, inner sep=5pt, align=center, font=\small, minimum height=3.5ex},
      %
      % Larghezze ottimizzate (13.5em per i centrali, 9.5em per i laterali)
      host/.style={cls, fill=orange!10, draw=orange!50, minimum width=13.5em},
      infra/.style={cls, fill=blue!7, draw=blue!40, minimum width=13.5em},
      %
      app/.style={cls, fill=green!7, draw=green!45, minimum width=9.5em},
      backend/.style={cls, fill=gray!8, draw=gray!50, minimum width=9.5em},
      users/.style={cls, fill=gray!8, draw=gray!50, minimum width=9.5em},
      %
      group/.style={draw=gray!60, fill=gray!3, dashed, inner sep=12pt, rounded corners=3pt},
      %
      % Frecce native
      arr/.style={-latex, semithick, gray!80!black},
      darr/.style={-latex, semithick, dashed, gray!80!black},
      barr/.style={latex-latex, semithick, gray!80!black},
    ]

    %% ── Asse Centrale ─────────────────────────────────────────────────────
    \node[host] (host) {\textbf{Host: apiguard-assurance}\\{\footnotesize Python process, \texttt{apiguard run}}};

    % Spazio verticale ridotto a 7ex per compattare
    \node[infra, below=7ex of host] (kong) {\textbf{Kong~3.9 DB-less}\\{\footnotesize :8000 (HTTP proxy)}\\{\footnotesize :8443 (HTTPS proxy)}\\{\footnotesize :8001 (Admin \gls{API})}};

    %% ── Ali Laterali (Distanza 2.5em per restare nei margini) ─────────────
    \node[app, right=2.5em of kong] (forgejo) {\textbf{Forgejo~14.0.3}\\{\footnotesize target applicativo}};

    \node[users, left=2.5em of kong] (rbac) {\textbf{Utenti RBAC}\\{\footnotesize thesis-admin}\\{\footnotesize user-a (std.)}\\{\footnotesize user-b (GREY\_BOX~\#2)}};

    %% ── Backend ───────────────────────────────────────────────────────────
    \node[backend, below=4ex of forgejo] (pg) {\textbf{PostgreSQL}\\{\footnotesize backend Forgejo}};

    %% ── Container Group (Docker Compose network) ──────────────────────────
    \begin{scope}[on background layer]
      \node[group, fit=(kong)(forgejo)(pg)] (compose) {};
    \end{scope}
    \node[font=\footnotesize\bfseries, text=gray!70!black, anchor=south west, xshift=1ex, yshift=0.5ex] at (compose.south west) {Docker Compose network};

    %% ── Frecce ────────────────────────────────────────────────────────────
    \draw[arr] (host.south) -- node[pos=0.35, right, align=left, font=\footnotesize\itshape, text=gray!70!black] {:8000/:8443 probe \\ :8001 Admin API} (kong.north);

    % Testo spezzato su due righe tramite align=center e \\
    \draw[arr] (kong.east) -- node[midway, above, align=center, font=\footnotesize\itshape, text=gray!70!black] {proxy/ \\route} (forgejo.west);

    \draw[barr] (forgejo.south) -- node[midway, right, font=\footnotesize\itshape, text=gray!70!black] {SQL} (pg.north);

    \draw[darr] (rbac.east) -- node[midway, above, font=\footnotesize\itshape, text=gray!70!black] {token} (kong.west);

  \end{tikzpicture}
  \caption{Topologia dell'ambiente di validazione sperimentale (Forgejo~14.0.3 + Kong~3.9 DB-less, Docker Compose).}
  \label{fig:testbed-topology}
\end{figure}
